%
% Documento: Resumo (Português)
%



\begin{resumo}

\vspace*{12pt}

\begin{SingleSpacing} 
            Esta pesquisa visa contribuir para a determinação de um parâmetro fundamental em astronomia: a massa de estrelas jovens em aglomerados estelares, com potencial para enriquecer a compreensão da rotação estelar. O objetivo é estimar massas estelares em aglomerados como o da Nebulosa de Órion (ONC), Plêiades, h Persei, Upper Scorpius e NGC 2516 utilizando uma combinação de métodos tradicionais, como ajuste de isócronas e relações massa-luminosidade, e técnicas de aprendizado de máquina. Este estudo adota uma abordagem descritiva para investigar métodos de cálculo de massa estelar, utilizando ferramentas estatísticas para avaliar a qualidade das estimativas. A amostra é não probabilística, com foco em estrelas jovens de baixa massa, e os dados foram coletados de diversos catálogos, incluindo o GAIA e o TESS. As magnitudes absolutas foram calculadas a partir da fotometria, e a extinção interestelar individual foi considerada para os aglomerados mais jovens. A massa estelar foi estimada utilizando dois métodos dependentes de modelos teóricos de evolução e interior estelar: o ajuste de isócronas e a relação massa-magnitude, esta utilizando técnicas de aprendizado de máquina, incluindo validação cruzada por K-fold e otimização por busca em grade, com diversas técnicas de regressão como Florestas Aleatórias, Regressão com Vetor de Suporte, e Regressão Bayesiana. Os resultados foram avaliados com um conjunto de métricas estatísticas, como RMSE, AIC e Coeficiente de determinação, comparando com valores de massa obtidos na literatura, demonstrando a eficácia da relação massa-luminosidade que teve uma boa concordância com a maioria dos aglomerados, principalmente os com uma sequência principal bem definida, como as Plêiades, e do ajuste de isócronas para aglomerados com mais objetos na pré-sequencia principal, como a ONC, embora com incertezas e limitações inerentes. A integração de técnicas de aprendizado de máquina se mostrou promissora para aumentar a precisão e confiabilidade das estimativas de massa estelar, apontando direções futuras para pesquisas astrofísicas.


         %   \textbf{Palavras-chave}: Parâmetros estelares: Massa. Modelagem Computacional. Estatística.

 
\end{SingleSpacing}
\end{resumo}
